\section{Discussion}
The reason for the western intensification can be explained by the phase movement of Rossby waves always being westward. Thus, larger fluxes appear on the western half. In the Munk theory, the imposed harmonic viscosity ensured that these Rossby waves died down, ensuring a stable solution. However, the lack of diffusion in the Sverdrup solution forced continuous Rossby wave propagation, which explains the non-converging result. One can equally explain this theoretically by a mass accumulation along the western boundary creating a high which propagates like a Rossby wave. This explains the clear limitation of the solution not reaching steady state, However, since the Rossby wave propagation was superimposed on the solution, time averaging resulted in good results. 

The Sverdrup non-convergence can be explained through the lack of energy dissipation. Through continuous energy input, propagation of Rossby waves was continuously fed. This is the same mechanism that forced the $f$-plane solution to grow continuously. Thus one needs some kind of dissipation to balance the continuous energy input. In the real world this is done through small-scale dissipation, connections to other oceans and other dynamical properties. 

Another problem of the Sverdrup solution is the lack of mass convergence in the theory. Since the theory is first order, only one boundary condition could be imposed, which cannot allow return flow (although this was forced numerically through mass conservation). The Munk solution is higher order, allowing multiple boundary conditions and a return flow. The return flow can be explained through the frictional effect opposing the mean southward flow, by creating a boundary layer close to the western boundary. Since frictional effects act in the opposite direction, a strong northern current is created. 

One can furthermore criticize the Munk model for artificial viscosity. If the Stommel drag is the friction caused by bottom drag, the Munk viscosity should be a boundary effect, and thus only be seen at the boundaries. However, although the viscosity was implemented throughout the entire domain, the Munk layer ensured that this effect was only present near the boundary. Furthermore, since the viscosity presents reasonable results, one can argue that the model is of good use. However, in most realistic models one often uses the combined Munk--Stommel viscosity to obtain benefits of both models \citep{Vallis06}. 

One can always obtain better results by increasing the grid size in the domain, for example to resolve the western part of the boundary current in the Munk model. Since the lowest viscosity only had a resolution of $\tfrac{\delta_M}{\Delta x}=3.8$, one can observe a non zero value at the western boundary. Thus increasing $\Delta x$ one could obtain a higher resolution and $\Psi\left(x_\text{West}\right)\rightarrow 0$. However, according to \citet{Steele2008eos}, the speed of the Gulf Stream is around \SI{2.5}{\meter\per\second}. Comparing this value with \autoref{fig:sverdrup_u10}--\ref{fig:munk_a}, the obtained results in these schemes are thus on the right order of magnitude compared to the real world. Thus, even though this article presented a simplified model, reasonable values were obtained. Furthermore, the theoretical error remained below 15\% for both models, showcasing working numerical models.